% ════════════════════════════════════════════
% 引言（01.引言.tex）—— 写作提示词见本文件末尾【本章节写作提示词】。
% ════════════════════════════════════════════
\section{引言}
\subsection{问题背景}

随着低空经济与无人机的快速发展，小型无人机在低空空域的违规飞行、侦察与袭扰对重要目标的安全防护构成严峻挑战。由于小型无人机雷达散射截面小、飞行高度低，传统雷达侦察系统往往只能给出包含大致方位角、水平距离和高度的粗略告警信息，无法直接得到目标的精确地理位置。因此，如何快速、准确地实现无人机目标的测向与定位，成为低空防御领域亟待解决的关键问题。

为实现快速测向与定位，指挥调度中心通常派出搭载云台相机的车载地面移动观测车前往有利位置对目标进行测向。单台观测车只能获得一条方位线，无法确定目标水平位置；两台观测车从不同方位同步观测，可通过两条方位线的交会实现定位。然而，两站之间的基线过短或交会角过大、过小都会放大定位误差，车辆还受到建筑物阻挡、视线遮挡和到达时间期限的约束，新的告警可能突然出现，车辆故障或通信中断也可能随时发生。这些因素使得站址选择、车辆指派与动态重调度成为一个复杂的组合优化问题。现有研究已对多站测向交叉定位及其精度分析开展了较多工作\cite{wang2007cexiang,liu2008cexiang}，为本研究提供了几何误差分析基础。

因此，研究多移动观测车协同的无人机快速定位问题，建立可观测区域判定、交会定位精度、误差传播与动态重调度的一体化数学模型，对于提升低空空域防御的快速响应能力具有重要的理论意义和工程价值。

\subsection{问题重述}

\textbf{问题1：} 在车辆几何、视线遮挡与云台俯仰三项约束下，确定目标 T1 的可观测区域，并以观测质量 Q 为准则确定最优观测站及其水平距离与俯仰角，给出并列取舍规则。

\textbf{问题2：} 构建两站交会定位精度模型，在行驶时间约束下确定告警 T2 的最优站对；并针对 T1、T2、T3 并发告警，设计车辆复用与路径规划方案以最小化最大观测开始时间，与最近指派启发式对比。

\textbf{问题3：} 构建车辆定位误差与云台指向误差对水平方位精度的误差传播模型，确定满足方位偏差不超过 5 度且概率不小于 0.95 的最大允许行驶距离，推导误差标准差模型中参数的约束，并比较三类误差源的相对贡献。

\textbf{问题4：} 以问题二所得调度方案为基线，针对新增告警、随机车辆故障与通信中断及其组合扰动建立重调度模型，与保持原方案策略比较，在 60 秒评价窗口内按完成率与最晚观测开始时间评估。
